\documentclass[a4paper,12pt]{article}

% Русский язык и кодировка
\usepackage[T2A]{fontenc}
\usepackage[utf8]{inputenc}
\usepackage[russian]{babel}

% Поля страницы
\usepackage[
    left=3cm,
    right=2cm,
    top=2cm,
    bottom=2cm
]{geometry}

% Работа с изображениями
\usepackage{graphicx}
\usepackage{float}

% Таблицы
\usepackage{array}
\usepackage{longtable}


% Ссылки
\usepackage[hidelinks]{hyperref}

% Межстрочный интервал
\usepackage{setspace}
\onehalfspacing

\begin{document}

% -------------------------------------------------
% Титульная страница
% -------------------------------------------------

\begin{titlepage}
    \centering

    {\large Университет ИТМО\par}
    \vspace{1cm}

    {\large Лабораторная работа №2\par}
    \vspace{0.5cm}

    {\Large \textbf{Документирование проекта и работа с Git}\par}

    \vfill

    \begin{flushright}
        Выполнил: Смирнов Иван Евгеньевич\\
        ИСУ: 561006
    \end{flushright}

    \vfill

    Санкт-Петербург\\
    2026
\end{titlepage}

% -------------------------------------------------
% Цель работы
% -------------------------------------------------

\section{Цель работы}

Целью лабораторной работы является изучение принципов документирования программного кода, создание документации в формате Markdown и работа с ветками и историей изменений в Git.

В рамках дополнительного задания также было выполнено оформление отчёта с использованием системы компьютерной вёрстки LaTeX.

% -------------------------------------------------
% Ход работы
% -------------------------------------------------

\section{Ход выполнения работы}

\subsection{Создание рабочей ветки}

Для выполнения лабораторной работы была создана отдельная ветка:

\begin{verbatim}
doc_561006
\end{verbatim}

Работа выполнялась локально в репозитории проекта, после чего изменения были опубликованы на GitHub.

\subsection{Анализ файла mysterious\_file.py}

В проекте был найден файл \texttt{mysterious\_file.py}, содержащий несколько функций с неинформативными названиями переменных и функций.

Для каждой функции были определены:
\begin{itemize}
    \item назначение;
    \item параметры;
    \item возвращаемое значение;
    \item более понятное название;
    \item документация в формате Python docstring.
\end{itemize}

В результате были получены следующие функции.

\begin{longtable}{|p{5cm}|p{9cm}|}
\hline
\textbf{Функция} & \textbf{Назначение} \\
\hline
\texttt{buy\_products} & Проверка возможности покупки и расчёт остатка баланса. \\
\hline
\texttt{is\_weekend} & Проверка того, является ли день недели выходным. \\
\hline
\texttt{is\_even} & Проверка числа на чётность. \\
\hline
\texttt{find\_multiples\_of\_three} & Поиск чисел, кратных трём, в заданном диапазоне. \\
\hline
\texttt{calculate\_sentiment\_score} & Расчёт оценки тональности текста. \\
\hline
\texttt{convert\_to\_binary} & Перевод числа в двоичное представление. \\
\hline
\end{longtable}

\subsection{Документирование функций}

Для каждой функции был добавлен документационный блок с описанием назначения, параметров и возвращаемого значения.

Например, функция проверки чётности получила следующий вид:

\begin{verbatim}
def is_even(number):
    """
    Проверяет, является ли число чётным.

    Args:
        number (int): число для проверки.

    Returns:
        bool: True, если число чётное, иначе False.
    """
    return number % 2 == 0
\end{verbatim}

Кроме того, неинформативные имена переменных были заменены на более понятные. При этом логика исходных функций была сохранена.

% -------------------------------------------------
% Markdown
% -------------------------------------------------

\subsection{Создание Markdown-документации}

В проекте был создан каталог \texttt{docs}, содержащий файл:

\begin{verbatim}
docs/functions.md
\end{verbatim}

В файле приведено описание всех функций проекта, их параметров, возвращаемых значений и примеров вызова.

Также в документации приведена история изменений с хешами Git-коммитов.

% -------------------------------------------------
% Git
% -------------------------------------------------

\section{История изменений Git}

В процессе выполнения работы были созданы осмысленные коммиты.

\begin{longtable}{|p{4cm}|p{9cm}|}
\hline
\textbf{Коммит} & \textbf{Описание} \\
\hline
\texttt{Document and rename functions} &
Документирование функций и переименование функций и переменных. \\
\hline
\texttt{Add function documentation} &
Добавление Markdown-документации проекта. \\
\hline
\texttt{Update documentation history} &
Обновление истории изменений в документации. \\
\hline
\end{longtable}

После завершения работы ветка \texttt{doc\_561006} была объединена с веткой:

\begin{verbatim}
main_features_561006
\end{verbatim}

Итоговая ветка была опубликована на GitHub.

% -------------------------------------------------
% Результат
% -------------------------------------------------

\section{Результат работы}

В результате лабораторной работы:

\begin{itemize}
    \item были документированы функции файла \texttt{mysterious\_file.py};
    \item функции и переменные получили информативные названия;
    \item был создан каталог \texttt{docs};
    \item была создана Markdown-документация проекта;
    \item была сохранена история изменений Git;
    \item рабочая ветка была объединена с \texttt{main\_features\_561006};
    \item итоговые изменения были опубликованы на GitHub.
\end{itemize}

% -------------------------------------------------
% Скриншоты
% -------------------------------------------------

\section{Скриншоты}

\subsection{Создание рабочей ветки}

Для выполнения лабораторной работы была создана отдельная ветка
\texttt{doc\_561006}.

\begin{figure}[H]
    \centering
    \includegraphics[width=0.9\textwidth]{image1.png}
    \caption{Создание рабочей ветки}
\end{figure}


\subsection{Изменение файла mysterious\_file.py}

В исходном файле функции имели неинформативные названия и содержали
необходимость в дополнительном документировании. В процессе работы
функции и переменные были переименованы, а для каждой функции добавлен
документационный блок.

\begin{figure}[H]
    \centering
    \includegraphics[width=0.9\textwidth]{image4.png}
    \caption{Документированный файл mysterious\_file.py}
\end{figure}

\begin{figure}[H]
    \centering
    \includegraphics[width=0.9\textwidth]{image5.png}
    \caption{Продолжение изменённого файла mysterious\_file.py}
\end{figure}

\begin{figure}[H]
    \centering
    \includegraphics[width=0.9\textwidth]{image6.png}
    \caption{Окончание изменённого файла mysterious\_file.py}
\end{figure}


\subsection{Markdown-документация}

В каталоге \texttt{docs} был создан файл \texttt{functions.md},
содержащий описание функций, их параметров, возвращаемых значений,
примеров вызова и истории изменений.

\begin{figure}[H]
    \centering
    \includegraphics[width=0.9\textwidth]{image7.png}
    \caption{Markdown-документация проекта}
\end{figure}

\begin{figure}[H]
    \centering
    \includegraphics[width=0.9\textwidth]{image8.png}
    \caption{Продолжение Markdown-документации}
\end{figure}


\subsection{История изменений Git}

В процессе выполнения работы были созданы осмысленные коммиты.
Историю изменений можно увидеть с помощью команды \texttt{git log}.

\begin{figure}[H]
    \centering
    \includegraphics[width=0.9\textwidth]{image9.png}
    \caption{История изменений Git}
\end{figure}


\subsection{Публикация рабочей ветки}

После выполнения основной части работы ветка \texttt{doc\_561006}
была опубликована на GitHub.

\begin{figure}[H]
    \centering
    \includegraphics[width=0.9\textwidth]{image11.png}
    \caption{Рабочая ветка doc\_561006 на GitHub}
\end{figure}


\subsection{Финальное слияние}

После завершения работы ветка \texttt{doc\_561006} была объединена
с веткой \texttt{main\_features\_561006}. Итоговая ветка также была
опубликована на GitHub.

\begin{figure}[H]
    \centering
    \includegraphics[width=0.9\textwidth]{image13.png}
    \caption{Итоговая ветка main\_features\_561006 на GitHub}
\end{figure}

\end{document}